\documentclass[11pt]{article}

% --- Visuals: Mimic OpenReview Web Style ---
\usepackage[margin=1in]{geometry}
\usepackage{xcolor}
\renewcommand{\familydefault}{\sfdefault} 

% --- THE FIXES ---
\usepackage{parskip} % Removes paragraph indentation and adds space between paragraphs
\setcounter{secnumdepth}{0} % Removes numbering from all sections/headings (e.g., 1.1)

% --- Markdown Package Setup ---
\usepackage[
    hybrid=true,          
    pipeTables=true,      
    fencedCode=true,      
    hashEnumerators=true, 
    smartEllipses=true,
    underscores=false     
]{markdown}

\begin{document}

\begin{markdown}
We thank the reviewer for the detailed reproducibility suggestions. We agree that these details should be explicit and will add them, together with the executable measurement pipeline and directly reported timing variation, to the revised manuscript and artifact.

**Reproducibility and baselines.** The primary measurements used an NVIDIA RTX A6000 (compute capability 8.6, 47.5 GiB), Python 3.12.12, PyTorch 2.6.0+cu124, CUDA 12.4, cuDNN 9.1.0, driver 550.163.01, nvcc 12.4.99, and GCC/G++ 11.4.0. The cross-hardware sweep used an NVIDIA RTX PRO 6000 Blackwell Max-Q (compute capability 12.0, 95.0 GiB), Python 3.12.12, PyTorch 2.9.1+cu130, CUDA 13.0, cuDNN 9.13.0, driver 595.71.05, nvcc 13.0 (build 36424714_0), and GCC/G++ 11.4.0. The reference is the official Finder et al. implementation at [commit 147a75b](https://github.com/BGU-CS-VIL/WTConv/commit/147a75b7aedc8c0bebcf8c0b489819bbf9a1ebe8). Layer benchmarks use eager execution for both implementations, while end-to-end benchmarks apply `torch.compile` uniformly to all models. We did not add a custom fusion path only to the reference; it uses the released PyTorch operators and normal cuDNN dispatch. The CUDA extension is JIT-compiled with `-O3 --use_fast_math -DHAAR_MAX_K=7`. We will document the build and table-generation commands, code layout, and runtimes; the primary layer sweep contains approximately 12 minutes of warmup and timed kernel work, excluding first-use JIT compilation.

**Timing variation.** Each configuration was collected in one benchmark sweep, with 20 warmup iterations followed by 50 individually timed CUDA-event iterations. Thus, these are within-run samples, not 50 independent runs. The revised absolute $k=5$ latency table reports mean $\pm$ standard deviation. Since each displayed latency is the geometric mean over nine tensor shapes, the accompanying deviation is likewise the geometric mean of the nine within-configuration standard deviations. Across all 504 A6000 measurements, the median coefficient of variation is 0.57%, and approximately 90% of configurations are below 1.66%.

**FP16 stability.** We thank the reviewer for raising this point. The reformulation is algebraically exact: in exact arithmetic, it defines an identical mathematical operator. The only numerical difference is the order in which floating-point reductions are evaluated. FP16 has 10 fraction bits, giving a spacing of approximately $9.8\times10^{-4}$ near unit magnitude; differences of a few rounding units are therefore expected when multiple operations are reassociated. The observed maximum absolute difference is $2.0\times10^{-3}$ for both outputs and input gradients and does not grow beyond this value through $L=5$. Its magnitude and bounded behavior are consistent with ordinary FP16 rounding rather than numerical instability. Since the paper makes no convergence or accuracy claim, we do not believe an additional end-to-end training experiment is required to support its claims.

**End-to-end protocol.** ConvNeXt-T and WTConvNeXt-T use stage depths $(3,3,9,3)$ and widths $(96,192,384,768)$; WTConvNeXt-T uses levels $(5,4,3,2)$ and $k=5$, while ConvNeXt-T uses $k=7$. We measure batch size 64 at $224\times224$ in FP32, with 20 warmup and 50 timed batches, and apply `torch.compile` to every model. Inputs and targets are already GPU-resident, so data loading and host-to-device transfer are excluded. Inference measures the forward pass. A training step includes gradient clearing, forward cross-entropy, backward propagation, and an SGD update, with CUDA synchronization around the timed region. We will make these inclusions and exclusions explicit in the revised table and appendix.

\end{markdown}

\end{document}
